%% Reference CV layout. Started from ai-job-search's moderncv example (MIT, see LICENSE-ai-job-search).
%% A starting point, not a form: write each CV for its job. Keep \usepackage{winnow-cv}: it sets the fonts
%% (Montserrat headings, Open Sans body), margins and spacing, and keeps contact icons out of the text layer.
%% Winnow compiles with Tectonic (XeTeX); winnow-cv.sty and fonts/ are copied next to the file first.

\documentclass[11pt,a4paper,sans]{moderncv}
\moderncvstyle{banking}
\moderncvcolor{blue}
\usepackage{winnow-cv}
\usepackage{needspace}

\AtEndPreamble{\hypersetup{pdftitle={[First] [Last] - CV}}}

% Personal data. Leave out any line you have nothing for (an empty \address{}{}{} fails to compile).
\name{[First]}{[Last]}
\phone[mobile]{[+XX XXXXXXXXXX]}
\email{[your.email@example.com]}
\extrainfo{[City, Country] \,|\, \href{[https://linkedin.com/in/you]}{[linkedin.com/in/you]} \,|\, \href{[https://github.com/you]}{[github.com/you]}}

\begin{document}

\makecvtitle

% Profile: three to five lines on why you fit this job, strongest match first.
[Profile statement.]

\section{Skills}
\begin{itemize}
\item \textbf{[Group 1]}: [Specific skills, tools, frameworks.]
\item \textbf{[Group 2]}: [Specific skills, tools, frameworks.]
\item \textbf{[Group 3]}: [Domain expertise, methods.]
\end{itemize}

\section{Experience}

% \cventry{dates}{title}{employer}{place}{}{details}. Dates use a single ASCII hyphen.
\cventry{[Mon YYYY - Present]}{[Job Title]}{[Company]}{[City, Country]}{}{%
\begin{itemize}
\item {[Achievement, with a number where the files have one.]}
\item {[Achievement.]}
\item {[Achievement.]}
\end{itemize}}

\cventry{[Mon YYYY - Mon YYYY]}{[Job Title]}{[Company]}{[City, Country]}{}{%
\begin{itemize}
\item {[Achievement.]}
\item {[Achievement.]}
\end{itemize}}

% An older role can be a single line.
\cventry{[Mon YYYY - Mon YYYY]}{[Job Title]}{[Company]}{[City, Country]}{}{}

\section{Education}

\cventry{[YYYY - YYYY]}{[Degree], [Field]}{[Institution]}{[City, Country]}{}{[Relevant modules or thesis.]}

% Optional sections, when they help this job: Languages, Publications, Projects, Awards.
\section{Languages}
\cvitem{}{[Language] (native), [Language] (fluent).}

\end{document}
